\documentclass[aspectratio=169, xcolor=dvipsnames]{beamer}

% --- Theme Setup ---
\usetheme{Madrid}
\usecolortheme{default}

% --- Packages ---
\usepackage{graphicx}
\usepackage{xcolor}
\usepackage{booktabs}
\usepackage{hyperref}
\usepackage{adjustbox}
\usepackage{amssymb}
\usepackage{amsmath}
\usepackage{listings}
\usepackage{caption}
\usepackage{tikz}

\lstset{
  basicstyle=\ttfamily\small,
  keywordstyle=\color{blue}\bfseries,
  stringstyle=\color{red},
  showstringspaces=false,
  breaklines=true
}

% --- Title Information ---
\title{Module 32}
\subtitle{Application Design and Development/2: Web Applications}
\author{Partha Pratim Das}
\institute{Department of Computer Science and Engineering\\Indian Institute of Technology, Kharagpur\\ppd@cse.iitkgp.ac.in}
\date{\today}

\begin{document}

% Title Slide
\begin{frame}
    \titlepage
\end{frame}

\begin{frame}{Module Recap}
    \begin{itemize}[<+->]
        \item Had a glimpse of Application Programs across various sectors
        \item Understood the typical architecture for an application
        \item Studies the classification and evolution of the architectures
        \item Glimpsed at architecture for a few sample applications
    \end{itemize}
\end{frame}

\begin{frame}{Module Objectives}
    \begin{block}{Goals}
        \begin{itemize}[<+->]
            \item To familiarize with the fundamentals notions and technologies of Web
            \item To learn about scripting
            \item To learn about Servlets
        \end{itemize}
    \end{block}
\end{frame}

\begin{frame}{Module Outline}
    \begin{itemize}[<+->]
        \item Web Fundamentals and Scripting
        \item Servlets
    \end{itemize}
\end{frame}

\section{Web Fundamentals}
\begin{frame}
  \vfill
  \centering
  \begin{beamercolorbox}[sep=8pt,center,shadow=true,rounded=true]{title}
    \usebeamerfont{title}Web Fundamentals\par%
  \end{beamercolorbox}
  \vfill
\end{frame}

\begin{frame}{The World Wide Web}
    \begin{itemize}[<+->]
        \item The Web is a distributed information system based on hypertext
        \item Most Web documents are hypertext documents formatted via the HyperText Markup Language (HTML)
        \item HTML documents contain
        \begin{itemize}
            \item[$\triangleright$] text along with font specifications, and other formatting instructions
            \item[$\triangleright$] hypertext links to other documents, which can be associated with regions of the text
            \item[$\triangleright$] forms, enabling users to enter data which can then be sent back to the Web server
        \end{itemize}
    \end{itemize}
\end{frame}

\begin{frame}{Uniform Resources Locators}
    \begin{itemize}[<+->]
        \item On the Web, functionality of pointers is provided by Uniform Resource Locators (URLs).
        \item \textbf{URL example:} \url{http://www.acm.org/sigmod}
        \begin{itemize}
            \item[$\triangleright$] The first part indicates how the document is to be accessed (protocol): 'http' indicates HTTP
            \item[$\triangleright$] The second part gives the unique name of a machine on the Internet
            \item[$\triangleright$] The rest of the URL identifies the document within the machine
        \end{itemize}
        \item The local identification can be:
        \begin{itemize}
            \item[$\triangleright$] The path name of a file on the machine
            \item[$\triangleright$] An identifier of a program plus arguments to be passed to the program
        \end{itemize}
    \end{itemize}
\end{frame}

\begin{frame}{URI, URL, and URN}
    \begin{block}{Definitions}
        \begin{itemize}[<+->]
            \item \textbf{Uniform Resource Identifier (URI)}: The overarching concept
            \item \textbf{Uniform Resource Locator (URL)}: Locates the resource
            \item \textbf{Uniform Resource Name (URN)}: Names the resource
        \end{itemize}
    \end{block}
    \begin{itemize}[<+->]
        \item URIs can be classified as locators (URLs), or as names (URNs), or as both
        \item URN functions like a person's name; URL resembles that person's street address
        \item URN defines an item's identity, while the URL provides a method for finding it
    \end{itemize}
\end{frame}

\begin{frame}{HTML and HTTP}
    \begin{itemize}[<+->]
        \item HTML provides formatting, hypertext link, and image display features
        \begin{itemize}
            \item[$\triangleright$] including tables, stylesheets (to alter default formatting), etc.
        \end{itemize}
        \item HTML also provides input features
        \begin{itemize}
            \item[$\triangleright$] Select from a set of options: Pop-up menus, radio buttons, check lists
            \item[$\triangleright$] Enter values: Text boxes
        \end{itemize}
        \item Filled in input sent back to the server, to be acted upon by an executable at the server
        \item HyperText Transfer Protocol (HTTP) used for communication with the Web server
    \end{itemize}
\end{frame}

\begin{frame}[fragile]{Sample HTML}
    \begin{verbatim}
<html>
<body>
<table border>
<tr><th>ID</th><th>Name</th><th>Department</th></tr>
<tr><td>00128</td><td>Zhang</td><td>Comp. Sci.</td></tr>
...
</table>
<form action='PersonQuery' method='get'>
Search for:
<select name='persontype'>
<option value='student' selected>Student</option>
<option value='instructor'>Instructor</option>
</select>
<br>
Name:
<input type='text' size='20' name='name'>
<input type='submit' value='submit'>
</form>
</body>
</html>
    \end{verbatim}
\end{frame}

\begin{frame}{HTTP and Sessions}
    \begin{itemize}[<+->]
        \item \textbf{The HTTP protocol is connectionless}
        \begin{itemize}
            \item[$\triangleright$] Once the server replies, it closes the connection and forgets the request
            \item[$\triangleright$] Unix logins and JDBC/ODBC connections stay connected, retaining authentication
            \item[$\triangleright$] \textbf{Motivation:} reduces load on server (OS has tight limits on open connections)
        \end{itemize}
        \item Information services need session information
        \begin{itemize}
            \item[$\triangleright$] For example, user authentication should be done only once per session
        \end{itemize}
        \item \textbf{Solution:} use a cookie
    \end{itemize}
\end{frame}

\begin{frame}{Sessions and Cookies}
    \begin{itemize}[<+->]
        \item A \textbf{cookie} is a small piece of text containing identifying information
        \begin{itemize}
            \item[$\triangleright$] Sent by server to browser on first interaction to identify session
            \item[$\triangleright$] Sent by browser back to the server that created the cookie on further interactions (part of HTTP protocol)
        \end{itemize}
        \item Server saves information about cookies it issued, and can use it when serving a request
        \begin{itemize}
            \item[$\triangleright$] For example, authentication information, and user preferences
        \end{itemize}
        \item Cookies can be stored permanently or for a limited time
    \end{itemize}
\end{frame}

\begin{frame}{Web Browser}
    \begin{itemize}[<+->]
        \item A web browser is application software for accessing the World Wide Web
        \item A web browser fetches content from the Web and displays it on a user's device
        \item A browser or rendering engine transforms HTML documents and other resources into an interactive visual representation
        \begin{itemize}
            \item[$\triangleright$] This includes image and video formats supported by the browser
        \end{itemize}
        \item Web pages usually contain hyperlinks to other pages and resources
        \item Web browsers are used on a range of devices, including desktops, laptops, tablets, and smartphones
        \begin{itemize}
            \item[$\triangleright$] In 2020, an estimated 4.9 billion people used a browser
            \item[$\triangleright$] Most used browser: Google Chrome with 64\% global market share; Safari with 19\%
        \end{itemize}
    \end{itemize}
\end{frame}

\begin{frame}{Web Servers}
    \begin{itemize}[<+->]
        \item A web server is software and underlying hardware that accepts requests via HTTP or HTTPS
        \item A web browser or crawler requests a specific resource using HTTP, and the server responds with the content or an error message
        \item The document name in a URL may identify an executable program that, when run, generates a HTML document
        \item To install a new service on the Web, one simply needs to create and install an executable that provides that service
        \item Common Gateway Interface (CGI): a standard interface between web and application server
    \end{itemize}
\end{frame}

\begin{frame}{Web Services}
    \begin{itemize}[<+->]
        \item Allow data on Web to be accessed using remote procedure call mechanism
        \item Two approaches are widely used
        \begin{itemize}
            \item[$\triangleright$] \textbf{REST (Representation State Transfer):} allows use of standard HTTP request to a URL to execute a request and return data; encoded in XML or JSON
            \item[$\triangleright$] \textbf{Big Web Services:} uses XML representation for sending request data and returning results; standard protocol layer built on top of HTTP
        \end{itemize}
    \end{itemize}
\end{frame}

\section{Scripting}
\begin{frame}
  \vfill
  \centering
  \begin{beamercolorbox}[sep=8pt,center,shadow=true,rounded=true]{title}
    \usebeamerfont{title}Scripting\par%
  \end{beamercolorbox}
  \vfill
\end{frame}

\begin{frame}{Two-Tier Web Architecture}
    \begin{center}
%         \includegraphics[height=0.6\textheight,keepaspectratio]{"../Lec 7.2_artifacts/image_000019_9c34e7bbcf12df880d5ddfe9b67425f171bc8a846c9a3dc52e697843d1a89c89.png"}
    \end{center}
\end{frame}

\begin{frame}{Three-Tier Web Architecture}
    \begin{center}
%         \includegraphics[height=0.6\textheight,keepaspectratio]{"../Lec 7.2_artifacts/image_000021_641e7d8c067e2a967ea9c42416b9d6dae31bf4e223edb39352e4ca06d9a16f73.png"}
    \end{center}
\end{frame}

\begin{frame}{Client-Side Scripting and Validation}
    \begin{itemize}[<+->]
        \item Advantages of execution at client
        \begin{itemize}
            \item[$\triangleright$] Offloads work from server
            \item[$\triangleright$] Provides more flexible user interface
        \end{itemize}
        \item \textbf{Disadvantages}
        \begin{itemize}
            \item[$\triangleright$] Client functionality is not guaranteed (e.g., Javascript may be disabled)
            \item[$\triangleright$] Browser may not support newer language features
            \item[$\triangleright$] Same application looks/behaves differently in different browsers
            \item[$\triangleright$] Client code can be hacked; security checks must always be verified at the server
        \end{itemize}
    \end{itemize}
\end{frame}

\begin{frame}{Javascript}
    \begin{itemize}
        \item \textbf{Javascript}
        \begin{itemize}
            \item Language very similar to Java
            \item Typically executed by the browser
        \end{itemize}
        \item Javascript can
        \begin{itemize}
            \item validate input,
            \item interact with user,
            \item alter display based on user interaction
        \end{itemize}
        \item Document Object Model (DOM)
        \begin{itemize}
            \item Tree representation of HTML document
            \item Nodes of tree can be accessed via Javascript
            \item $\triangleright$ \texttt{document.getElementById('some\_id')}
            \item And updated
            \item $\triangleright$ \texttt{document.getElementById('some\_id').innerHTML = "<b>Hello</b>"}
            \item Forms can be checked, and values updated
        \end{itemize}
    \end{itemize}
\end{frame}

\begin{frame}[fragile]{Javascript (2): Example}
    \begin{verbatim}
<html>
<head>
<script type="text/javascript">
function validate() {
  var creds = document.getElementById("CREDIT_CARD").value;
  if (creds.length != 16) {
    alert("Credit card must have 16 digits");
    return false;
  } else {
    return true;
  }
}
</script>
</head>
<body>
<form action="form_action.asp" onsubmit="return validate()">
  Name: <input type="text" name="name" size="20"> <br>
  Credit Card: <input type="text" name="CREDIT_CARD" id="CREDIT_CARD" size="16">
  <input type="submit" value="Submit">
</form>
</body>
</html>
    \end{verbatim}
\end{frame}

\begin{frame}{Server-Side Scripting}
    \begin{itemize}[<+->]
        \item Server-side scripting simplifies the task of connecting a database to the Web
        \begin{itemize}
            \item[$\triangleright$] Define an HTML document with embedded executable code/SQL queries
            \item[$\triangleright$] Input values from HTML forms can be used directly in the embedded code/SQL queries
            \item[$\triangleright$] When the document is requested, the Web server executes the embedded code/SQL queries to generate the actual HTML document, which is sent back to the client
        \end{itemize}
    \end{itemize}
\end{frame}

\begin{frame}{Server Side Context}
    \begin{itemize}[<+->]
        \item \textbf{A Web Application Layer}
        \begin{itemize}
            \item[$\triangleright$] Java Servlets
            \item[$\triangleright$] Java Server Pages (JSP)
            \item[$\triangleright$] Application Programming Interface (API) -- Java code with annotations
            \item[$\triangleright$] Python
            \item[$\triangleright$] Node JS
        \end{itemize}
        \item Web application layer maps HTTP requests to method invocations on a class
        \begin{itemize}
            \item[$\triangleright$] Maps parameters from the HTTP request to arguments of the method
            \item[$\triangleright$] Expects method to return a string or an object mapped to a string (such as HTML or JSON)
        \end{itemize}
    \end{itemize}
\end{frame}

\section{Servlets}
\begin{frame}
  \vfill
  \centering
  \begin{beamercolorbox}[sep=8pt,center,shadow=true,rounded=true]{title}
    \usebeamerfont{title}Servlets\par%
  \end{beamercolorbox}
  \vfill
\end{frame}

\begin{frame}{Servlets}
    \begin{itemize}[<+->]
        \item Java classes that run in a Web server
        \item Provide an API for HTTP interaction
        \begin{itemize}
            \item[$\triangleright$] Can process HTTP/HTTPS requests
            \item[$\triangleright$] Obtain parameters
            \item[$\triangleright$] Return results to the HTTP client (usually a web browser)
            \item[$\triangleright$] Maintain information across requests (Sessions)
        \end{itemize}
        \item \textbf{Example:} \texttt{javax.servlet.http.HttpServlet}
        \begin{itemize}
            \item[$\triangleright$] \texttt{doGet(request, response)}
            \item[$\triangleright$] \texttt{doPost(request, response)}
        \end{itemize}
        \item HTTP Methods
        \begin{itemize}
            \item[$\triangleright$] \textbf{GET}: pass parameters in URL
            \item[$\triangleright$] \textbf{POST}: pass parameters in the body of HTTP request
        \end{itemize}
    \end{itemize}
\end{frame}

\begin{frame}{Servlet (2)}
    \begin{center}
%         \includegraphics[height=0.7\textheight,keepaspectratio]{"../Lec 7.2_artifacts/image_000030_e06abfbbf6a04aacf1f8e120803c43105fd9c1cd5cd7ccb8accf58ef57cb3ef8.png"}
    \end{center}
\end{frame}

\begin{frame}[fragile]{Servlet (3): Example}
    \begin{verbatim}
String persontype = request.getParameter("persontype");
String number = request.getParameter("name");
if (persontype.equals("student")) {
  // ... code to find students with the specified name ...
  // ... using JDBC to communicate with the database ...
  out.println("<table BORDER COLS=3>");
  out.println("<tr><td>ID</td><td>Name</td><td>Department</td></tr>");
  for (/* each result */) {
    // ... retrieve ID, name and dept name
    // ... into variables ID, name and deptname
    out.println("<tr><td>" + ID + "</td><td>" + name + "</td><td>" + deptname + "</td></tr>");
  }
  out.println("</table>");
} else {
  // ... as above, but for instructors ...
}
    \end{verbatim}
\end{frame}

\begin{frame}{Servlet (4): Sessions}
    \begin{itemize}[<+->]
        \item Servlet API supports handling of sessions
        \item Sets a cookie on first interaction with browser, and uses it to identify session on further interactions
        \item To check if session is already active:
        \begin{itemize}
            \item[$\triangleright$] \texttt{if (request.getSession(false) == true)} $\Rightarrow$ existing session, else redirect to authentication
            \item[$\triangleright$] \texttt{request.getSession(true)}: creates new session
        \end{itemize}
        \item Store/retrieve attribute value pairs for a particular session
        \begin{itemize}
            \item[$\triangleright$] \texttt{session.setAttribute("userid", userid)}
            \item[$\triangleright$] \texttt{session.getAttribute("userid")}
        \end{itemize}
    \end{itemize}
\end{frame}

\begin{frame}{Servlet (5): Support}
    \begin{itemize}[<+->]
        \item Servlets run inside application servers such as
        \begin{itemize}
            \item[$\triangleright$] Apache Tomcat, Glassfish, JBoss
            \item[$\triangleright$] BEA Weblogic, IBM WebSphere and Oracle Application Servers
        \end{itemize}
        \item Application servers support
        \begin{itemize}
            \item[$\triangleright$] deployment and monitoring of servlets
            \item[$\triangleright$] Java 2 Enterprise Edition (J2EE) platform supporting objects, parallel processing across multiple application servers, etc
        \end{itemize}
    \end{itemize}
\end{frame}

\begin{frame}[fragile]{Java Server Pages (JSP)}
    \begin{itemize}
        \item A JSP page with embedded Java code
    \end{itemize}
    \begin{verbatim}
<html>
<head><title>Hello</title></head>
<body>
<% if (request.getParameter("name") == null) {
     out.println("Hello World");
   } else {
     out.println("Hello, " + request.getParameter("name"));
   } %>
</body>
</html>
    \end{verbatim}
    \begin{itemize}
        \item JSP is compiled into Java + Servlets
        \item JSP allows new tags to be defined, in tag libraries
        \begin{itemize}
            \item such tags are like library functions, can are used for example to build rich user interfaces such as paginated display of large datasets
        \end{itemize}
    \end{itemize}
\end{frame}

\begin{frame}[fragile]{PHP}
    \begin{itemize}
        \item PHP is widely used for Web server scripting
        \item Extensive libraries including for database access using ODBC
    \end{itemize}
    \begin{verbatim}
<html>
<head><title>Hello</title></head>
<body>
<?php
  if (!isset($_REQUEST['name'])) {
    echo "Hello World";
  } else {
    echo "Hello, " . $_REQUEST['name'];
  }
?>
</body>
</html>
    \end{verbatim}
\end{frame}

\begin{frame}{USP of JSP}
    \begin{itemize}[<+->]
        \item \textbf{JSP vs. Active Server Pages (ASP)}
        \begin{itemize}
            \item[$\triangleright$] ASP is a similar technology from Microsoft and is proprietary (uses VB)
            \item[$\triangleright$] JSP is platform independent and portable
        \end{itemize}
        \item \textbf{JSP vs. Pure Servlets}
        \begin{itemize}
            \item[$\triangleright$] JSP is more convenient to write and modify regular HTML than using millions of println statements
            \item[$\triangleright$] Web page designers can build the HTML, leaving places for servlet programmers to insert dynamic content
        \end{itemize}
        \item \textbf{JSP vs. JavaScript}
        \begin{itemize}
            \item[$\triangleright$] JavaScript can generate HTML dynamically on the client (Client Side)
            \item[$\triangleright$] JSP are executed by the web server before HTTP response (Server Side) -- can access server-side resources like databases
        \end{itemize}
        \item \textbf{JSP vs. Static HTML}
        \begin{itemize}
            \item[$\triangleright$] Regular HTML cannot contain dynamic information
        \end{itemize}
    \end{itemize}
\end{frame}

\begin{frame}{Module Summary}
    \begin{block}{Summary}
        \begin{itemize}[<+->]
            \item Familiarized with the Fundamentals notions and technologies of Web
            \item Learnt about Scripting
            \item Learnt the notions of Servlets
        \end{itemize}
    \end{block}
    \vspace{2em}
    \begin{center}
    \small \textit{Slides used in this presentation are borrowed from \url{http://db-book.com/} with kind permission of the authors.} \\
    \small \textit{Edited and new slides are marked with 'PPD'.}
    \end{center}
\end{frame}

\end{document}

